\section{Integrating the worker into Lean and Forge}
\label{sec:integration}
\status{The module layout, source bridge, routing rules, and acceptance gates
below are proposed. Only the Python components listed in Section~\ref{sec:evidence}
were executed. The six small Lean reconstruction targets are explicitly uncompiled.}

\subsection{Reuse the host architecture, not its scalar representation}

Forge's existing review of Leant and Djex already covers exact source ownership,
status distinctions, opaque receipts, and the separation of discovery from
verification~\cite{forge-leant}. Repeating that architecture as a contribution
would add little. Here its important consequence is concrete: a noncommutative
atom must retain its actual multiplication instance, scalar action, adjoint,
and matrix index type.

Leant's documented candidates retain source-level selections and are checked
again by Lean; Djex distinguishes complete propositional search from bounded
higher-rank synthesis~\cite{leant,djex}. The useful transfer is to treat an
operator expression proposed by those systems as a typed \emph{candidate}, not
as an algebraic theorem. Its semantics must still be established by the new
receipt or an existing proof. Neither project was built or invoked in these
experiments, and their historical test counts are not added to this report's
numbers.

The implementation should begin as an isolated \code{Forge.NC} worker at the
existing tactic/obligation boundary. It need not reach into \code{grind}'s
internal solver interfaces to test whether its new certificates help. Once a
reconstructed equality is available, return an ordinary Lean theorem to the
host's existing assertion mechanism. Operator and trace results should likewise
return the exact original proposition, not an untyped Boolean success flag.

\subsection{A proposed module split with narrow responsibilities}

\begin{longtable}{@{}p{0.29\textwidth}p{0.64\textwidth}@{}}
\toprule Module & Required responsibility \\\midrule\endhead
\code{Forge.NC.WordPoly} & Canonical finite rational word polynomials; addition,
concatenation, star, and evaluation theorems. No search state. \\
\code{Forge.NC.Reify} & Extract a supported expression and a proof that evaluating
the object polynomial gives the original Lean expression. Preserve atom types
and operation instances. \\
\code{Forge.NC.Receipt} & Typed equality, operator, trace, and matrix-witness
formats. Source problem is a separate argument. \\
\code{Forge.NC.Check} & Executable identity checks and their soundness theorems.
Use theorem-producing normalization first; reflection can replace it later. \\
\code{Forge.NC.Graded} & Optional formal account of homogeneous degree completeness
and finite quotient witnesses. Not required for accepting a positive equality
receipt. \\
\code{Forge.NC.Matrix} & Matrix evaluation, adjoint correspondence, congruence
positivity, and trace-specific reconstruction. \\
\code{Forge.NC.Search} & Untrusted context enumeration, slicing, Gram factorization,
ray proposal, and external backend adapters. \\
\code{Forge.NC.Worker} & Route supported obligations, enforce resource limits, and
return either reconstructed proofs or appropriately scoped diagnostics. \\
\bottomrule
\end{longtable}

A practical entrance vocabulary is:
\begin{lstlisting}
needEquality(sourceGoal, sourceRelations, algebraProfile)
needOperatorPositive(sourceGoal, sourceRelations, positiveFacts, starProfile)
needTracePositive(sourceGoal, sourceRelations, positiveFacts, traceProfile)
needAlgebraCountermodel(objectProblem, homogeneousEvidence, dimensionBudget)
\end{lstlisting}
These are interface sketches, not names of functions already present in Forge.
The last request must not receive authority to close an arbitrary Lean goal.

\subsection{A proof-producing source bridge}

For an expression $e$ over an algebra $A$, reification should return an object
polynomial $p$, an atom environment $\rho:X\to A$, and a Lean proof
\[
 \operatorname{eval}_{\rho}(p)=e.
\]
For every relation, it must return the analogous equation together with the
actual local proof that the source relation is zero. The atom map is indexed
by elaborated expressions in the exact context, not by display names.

The simplest supported grammar is literals, atoms, addition, subtraction,
negation, multiplication, natural powers, and a certified central rational
scalar action. For a nonrecognized subexpression, two safe choices exist:
retain it as an opaque atom of the same algebra, or decline the problem. An
opaque inverse is not supplied an inverse law automatically. For the
push-through example, the source must prove or explicitly supply the two inverse
relations; the worker never invents them from notation.

The bridge should establish its evaluation equations recursively. Atoms and
constants are immediate; the add and multiply cases use the corresponding
evaluation lemmas; a power follows by natural induction or a proved evaluation
homomorphism. Star uses the declared involution agreement for every atom. This
turns the search's word identity into a theorem about the original expression
without trusting serialization or pretty-printing.

Mathlib already provides \code{FreeAlgebra}, its universal lifting interface,
and an equivalence with a monoid algebra on the free monoid~\cite{free-algebra}.
These are useful semantic foundations. For executable checking, a canonical
array/list representation is likely more convenient than reducing a quotient
implementation directly. Prove its evaluation correspondence to the existing
semantic object once; do not formalize a second abstract free-algebra theory
merely to run sparse arithmetic.

\paragraph{Typed rectangular products.}
The delivered word language has one object: all atoms multiply in the same
algebra. Rectangular matrix products need typed words with a source and target
index type, and concatenation only when the middle types match. That extension
can use paths in a finite typed graph, but it is not part of the current
implementation. A preliminary adapter should reject mixed shapes rather than
pad matrices or forget indices. A trace rotation for rectangular factors also
moves between different matrix spaces; it needs the actual typed trace theorem,
not a same-size word permutation.

\subsection{Two reconstruction stages, with measurable gates}

\paragraph{Stage one: theorem-producing replay.}
Decode a small receipt, reconstruct each word as a source expression, and ask
existing normalization to prove the advertised free-algebra identity. For
integer-coefficient identities, \code{noncomm\_ring} is an immediate baseline
reconstructor. Prove the zero contribution of each relation by congruence,
multiplication by zero, and the supplied hypothesis. Sum the resulting
proofs. This route avoids formalizing the search or a Gr\"obner algorithm, while
allowing real kernel-checked examples early.

Rational coefficients need explicit scalar-action normalization. A tempting
shortcut is multiplying through by a common denominator. That also creates a
cancellation obligation: cancellation of a positive integer scalar is not valid
in every ring. On a $\Q$-algebra the scalar is invertible; the bridge must use
that structure explicitly. On an integer-only ring receipt, retain integer
coefficients and use the simpler reconstruction instead.

\paragraph{Stage two: reflected replay.}
When proof terms from expanded identities become the bottleneck, define a
computational checker over canonical word polynomials and prove
\[
 \operatorname{checkEq}(p,r,c)=\mathrm{true}
 \quad\Longrightarrow\quad
 \bigl(\forall j,\operatorname{eval}(r_j)=0\bigr)
 \Rightarrow\operatorname{eval}(p)=0.
\]
Analogous theorems interpret operator and trace receipts. The implementation
must connect checked object evaluation to source reification before assigning a
metavariable. A proved checker theorem still does not justify a mistranslated
source goal.

The supplied \code{lean/CertificatePrinciples.lean} contains six small targets
for this path: a sandwich-zero lemma, a finite-sum certificate principle, a
commuted square, the resolvent and push-through identities, and positivity under
matrix congruence. It does not implement the computational checker, certify all
119 JSON objects, or expose a new tactic. The compiler gate records
\code{NOT\_RUN} because no Lean executable was available.

\subsection{Concrete matrix library dependencies}

The matrix reconstruction should target \code{Matrix.PosSemidef} directly,
avoiding an ambiguous generic \code{0 <= M}. Mathlib documents the needed
operations: \code{Matrix.PosSemidef.conjTranspose\_mul\_mul\_same},\par
\code{Matrix.PosSemidef.add}, \code{Matrix.PosSemidef.smul},\par
\code{Matrix.posSemidef\_sum}, and
\code{Matrix.PosSemidef.trace\_nonneg}~\cite{matrix-posdef}.
The Loewner-order API can be introduced in the surface layer after confirming
its scoped instance; the certificate core need not depend on that notation.

For a sandwich receipt, instantiate the congruence theorem at the decoded $Q$
and the actual proof of $G$'s positive semidefiniteness. Use exact rational sign
proofs for $d_k\ge0$, scalar multiplication, and finite sums. The final word
identity replaces the positive sum by the source target. For trace receipts,
apply trace linearity and a proved cyclicity identity to the decoded
commutators; the resulting theorem is scalar-valued and stays scalar-valued
when returned to the host.

\subsection{Concrete optional search backends}

\begin{longtable}{@{}p{0.22\textwidth}p{0.32\textwidth}p{0.38\textwidth}@{}}
\toprule Backend & Useful existing facility & Adapter obligation and limitation \\\midrule\endhead
Python standard library & Exact fractions, dictionaries, deterministic worklists &
The only executed producer/replayer substrate in this package. Suitable as a
reference, not a production performance claim. \\
FLINT / python-flint & Rational matrix row reduction and solving, including
\code{fmpq\_mat.rref()} and \code{solve()}~\cite{flint} &
Convert exact numerators/denominators, retain column labels, and re-expand returned
solutions. Backend matrix conventions do not define receipt semantics. \\
NCAlgebra & \code{NCPoly}, \code{NCGBX}, and bounded
\code{NCPolyGroebner}; symbolic/numeric SDP support through
\code{NCSDP}~\cite{ncalgebra} &
Useful untrusted Wolfram Language search. Export original two-sided context
multipliers, not just a zero remainder. Requires a separate installed host;
not invoked here. \\
Ncpol2sdpa & Sparse moment/SDP relaxations for noncommutative polynomial
optimization~\cite{ncpol2sdpa,wittek} &
Use for candidate Gram data or supports. A moment relaxation bound or floating
solver verdict is not an accepted certificate. The appropriate primal/dual
conversion must be implemented and checked. \\
Mathlib & Free-algebra semantics, existing normalization, matrix positivity
lemmas~\cite{free-algebra,noncomm-ring,matrix-posdef} &
Proof reconstruction and final authority. Neither an installed compatible
package set nor a successful build is claimed by this dependency table. \\
\bottomrule
\end{longtable}

For mixed-degree SDP search, let $W$ contain the chosen words, introduce a
symmetric Gram matrix for each positive premise, and impose coefficient
equalities modulo the equality span. A numerical solve suggests supports or
approximate matrices. Reconstruct rational entries by solving the exact affine
coefficient equations on that support; then test each rational matrix by exact
$LDL$. If repair fails, reject the proposal. Merely rounding entries and checking
that a numerical eigenvalue is ``almost nonnegative'' is not a repair algorithm.

This is a proposed backend path. It extends the coefficient language in the
same way as the executed finite-ray search, but does not inherit that routine's
finite support completeness claim. Exact rational repair on a chosen numerical
support can miss feasible solutions. It must remain an untrusted candidate
producer with a distinct status.

For NCAlgebra, the current documentation recommends the pure-Wolfram
\code{NCGBX} route rather than the old C++ component~\cite{ncalgebra}. Even a
successful noncommutative completion is only a discovery mechanism here; a
nonterminating or iteration-limited completion yields \unknown{}. The
homogeneous degree procedure supplies a simpler exact alternative when its
entrance conditions hold.

\subsection{A routing algorithm the planner can implement}

\begin{lstlisting}
route(problem):
    reject unsupported scalar domain, source operations, or shape erasure
    try cheap existing normalization and proved local rewrites first
    if algebra equality:
        if all relations homogeneous:
            run context search at target degree
            replay identity if found
            otherwise construct and replay quotient matrices within budget
        else:
            run bounded context levels; all failures stay UNKNOWN
    if operator positivity:
        discharge or expose missing source self-adjointness obligations
        if free homogeneous even-degree target and no needed constraints:
            try unique-Gram exact LDL
        try relation-derived rays plus explicit source-supplied rays
        optionally request an external candidate Gram certificate
    if trace positivity:
        keep the trace wrapper
        use cyclic residual search in the implemented fragment
    export only a proof of the original proposition, or a scoped diagnostic
\end{lstlisting}

The homogeneous Gram fast path can safely ignore additional assumptions when it
proves the target unconditionally, but a failure on that weaker unconstrained
problem must not refute the constrained goal. Similarly, equality receipts may
use a subset of available relations for positive proof; algebra countermodels
must satisfy every relation included in the claimed universal object problem.
The original source context may be richer still, which is why a source-level
negative verdict needs a stronger translation gate.

A useful exchange with Leant-style typed synthesis is to request candidate
operator expressions $q$ from an exact local inventory, then pass those
expressions to the positive-ray search. A useful exchange in the other direction
is to return an equality that simplifies the behavioral specification of a
candidate. The first implementation need not support unrestricted higher-rank
synthesis: the executed involution dictionary already exercises the
proposal/verification split with fully explicit algorithms.

\subsection{Proof size and cache discipline specific to this worker}

Equality elimination provenance can become large even when the final target
has a short proof. Prefer a shared directed acyclic graph for repeated context
multiplications, followed by dead-term removal and exact coefficient collection.
The current JSON format expands the final list and does not implement DAG
compression. The slicing experiment shows a different benefit: a three-term
receipt can be found without building an elimination basis of rank 1,345.

Memoizing the equality remainder basis is useful across repeated demands from
one presentation. Its cache key must include the exact ordered atom meanings,
coefficient domain, complete relation list, context bound, and involution when
used for positivity. A cache may provide \emph{search guidance} after a looser
match; it may not provide an accepted proof. The graph dictionary can be cached
separately as proposals, but a cached square still needs its source typing and
current receipt replay.

The supplied checker imposes finite input sizes and product counts, rejects
booleans as integers, requires canonical rational encodings, rejects duplicate
JSON keys, and fails closed on malformed evidence. These are useful research
checks, not a claim of a hardened parser for hostile internet workloads. A
production process boundary also needs byte limits before JSON decoding,
cooperative cancellation, wall-clock and memory limits, and failure isolation.

\section{An incremental implementation plan with acceptance tests}
\label{sec:roadmap}

\subsection{Gate A: one real source-level equality}

Implement rational-word reification and theorem-producing equality replay.
Accept a Lean goal only after producing a proof of its exact source expression.
Run the resolvent and push-through examples in a fresh Lean process, record
compiler/toolchain and Mathlib revisions, and inspect the actual axiom
inventories. Include changed-target, changed-multiplication-instance, and
characteristic-two controls. The delivered prototype is not a substitute for
this gate.

\subsection{Gate B: reusable equality search}

Add exact context enumeration and incidence slicing. Compare the two paths on
identical Lean goals with the same reconstruction and resource limits. Measure
search time, reconstruction time, kernel time, peak memory, and proof size
separately. A small Python elimination time is not a prediction of Lean
elaboration time. Add multiple demands sharing a relation basis, but preserve
source-level context identities between demands.

\subsection{Gate C: operator and trace proof reconstruction}

Port the star correspondence and sandwich/trace interpretation theorems.
Require the self-adjoint projection and positive congruence examples to close
source goals. Then replay the homogeneous degree-six result and the two-stage
CHSH derivation. Reject the trace/operator confusion example as an operator
proof. No external SDP integration should precede these positive and negative
controls: the hard boundary is reconstruction, not sophisticated discovery.

\subsection{Gate D: exact-fragment diagnostics}

The first production countermodel feature can be a diagnostic for an explicitly
reified universal-algebra query. Verify whole relation matrices and the target
vector in Lean or a separately verified matrix checker. Only then consider
connecting it to source-goal refutation, with fixed dimensions and all local
assumptions accounted for. A proved impossibility result for the object language
must never become a proof of arbitrary Lean non-inhabitation.

\subsection{Gate E: broader search and a fair comparative corpus}

Only after the earlier gates should the project integrate a numerical SDP or
noncommutative completion backend. Construct a held-out corpus of symbolic
matrix, endomorphism, projection, and trace goals. Include failed goals and true
goals outside the worker's fragments, not just generated certificates.
Compare existing tactics, a noncommunicating portfolio, and the integrated
worker. Report solved original goals under a fixed budget, not counts of
internal algebra subgoals or generated receipt terms. The intended milestone
is a measured increase in kernel-checked source goals, not a larger Python
certificate corpus.
